% ChatGPT Directions 0 :  
% This is a Tbox Problem set for the following standards: 5.NBT.A.1, 5.NBT.A.2 
%--------------------------------------------------
\documentclass[12pt]{article}
\usepackage[a4paper, top=0.8in, bottom=0.7in, left=0.8in, right=0.8in]{geometry}
\usepackage{amsmath}
\usepackage{amsfonts}
\usepackage{latexsym}
\usepackage{graphicx}
\usepackage{fancyhdr}
\usepackage{tcolorbox}
\usepackage{enumitem}
\usepackage{setspace}
\usepackage[defaultfam,tabular,lining]{montserrat} % Font settings for Montserrat

% General Comment: Template for creating problem sets in a structured format with headers, titles, and sections.
% This document uses Montserrat font and consistent styles for exercises, problems, and performance tasks.

% -------------------------------------------------------------------
% Directions for LaTeX Styling and Content
% 1. Include a header with standards and topic title: \fancyhead[L]{\textbf{<Standards>: <Topic Title>}}.
% 2. Section Breakdown:
%    - Learning Objective: Concise goal statement.
%    - Exercises: Procedural fluency tasks.
%    - Problems: Moderately complex scenarios.
%    - Performance Task: Real-world, multi-step tasks.
%    - Reflection: Prompt to reflect on strategies and learning.
% 3. Styling with tcolorbox:
%    - Frame color: colframe=black!60.
%    - Background color: colback=gray!5 or white.
%    - Title Background: colbacktitle=black!15.
%    - Font Style: Bold and large (fonttitle=\bfseries\Large).
% -------------------------------------------------------------------

\setlength{\parindent}{0pt}
\pagestyle{fancy}

\setlength{\headheight}{27.11148pt}
\addtolength{\topmargin}{-15.11148pt}

\fancyhf{}
%\fancyhead[L]{\textbf{5.NBT.A.1, 5.NBT.A.2: Place Value and Patterns in Powers of 10}}
\fancyhead[R]{\includegraphics[width=0.8cm]{Round Logo.png}} % Placeholder for logo
\fancyfoot[C]{\footnotesize © Study Smart Tutors}

\sloppy

\title{}
\date{}
\hyphenpenalty=10000
\exhyphenpenalty=10000

\begin{document}

\subsection*{Problem Set: Place Value and Patterns in Powers of 10}
\onehalfspacing

% Learning Objective Box
\begin{tcolorbox}[colframe=black!40, colback=gray!5, 
coltitle=black, colbacktitle=black!20, fonttitle=\bfseries\Large, 
title=Learning Objective, halign title=center, left=5pt, right=5pt, top=5pt, bottom=15pt]
\textbf{Objective:} Understand place value to the thousandths and explain patterns when multiplying or dividing by powers of 10.
\end{tcolorbox}

% Exercises Box
\begin{tcolorbox}[colframe=black!60, colback=white, 
coltitle=black, colbacktitle=black!15, fonttitle=\bfseries\Large, 
title=Exercises, halign title=center, left=10pt, right=10pt, top=10pt, bottom=60pt]
\begin{enumerate}[itemsep=4em]
    \item Round \( 3,456.789 \) to the nearest whole number, tenth, and hundredth.
    \item What is the value of the digit \( 7 \) in \( 7,845 \)? What about in \( 7.845 \)?
    \item Multiply \( 3.45 \) by \( 10, 100, \) and \( 1,000 \). Explain the pattern in the products.
    \item Divide \( 5,600 \) by \( 10, 100, \) and \( 1,000 \). What happens to the decimal point in each case?
    \item Solve: \( (5,000 - 3,278) \div 2 \).
    \item Write and solve the equation: "A package weighs 144 ounces. If divided equally into 12 smaller boxes, how much does each box weigh?"
    \item Add and round: \( 12,345 + 45,678 \), rounded to the nearest thousand.
    \item Multiply \( 1,250 \times 8 \).
    \vspace{1cm}
\end{enumerate}
\end{tcolorbox}

\vspace{1em}

% Problems Box
\begin{tcolorbox}[colframe=black!60, colback=white, 
coltitle=black, colbacktitle=black!15, fonttitle=\bfseries\Large, 
title=Problems, halign title=center, left=10pt, right=10pt, top=10pt, bottom=100pt]
\begin{enumerate}[start=9, itemsep=5em]
    \item A factory produces \( 2,500 \) items in a day. How many items are produced in \( 10, 100, \) and \( 1,000 \) days? Explain the pattern in your results.
    \item A company packs \( 24,352 \) boxes in 4 days. How many boxes are packed per day? Write an equation to represent the scenario.
   
    \item A bakery sells \( 432 \) cookies in the morning and \( 576 \) cookies in the afternoon. If each box holds 24 cookies, how many boxes are used in total?
    \item A school orders \( 15,675 \) pencils. They distribute them equally to \( 25 \) classrooms. How many pencils does each classroom receive?
    \item Solve for \( x \): \( 8x + 500 = 1,100 \).
    \item Why does dividing \( 7,200 \) by \( 10 \) shift the digits one place to the right? Explain.
\end{enumerate}
\end{tcolorbox}

\vspace{1em}

% Performance Task Box
\begin{tcolorbox}[colframe=black!60, colback=white, 
coltitle=black, colbacktitle=black!15, fonttitle=\bfseries\Large, 
title=Performance Task: Designing a Playground, halign title=center, left=10pt, right=10pt, top=10pt, bottom=50pt]
You are designing a new playground. Here’s what you know:
\begin{itemize}
    \item The total area of the playground is 6,400 square feet.
    \item Half of the area is for a grassy field, and the other half is divided equally between a sandbox and a swing set.
    \item Each square foot of the playground costs \$10 to install.
\end{itemize}
\textbf{Task:}
\begin{enumerate}[itemsep=4em]
    \item Calculate the area of the grassy field.
    \item Calculate the area of the sandbox and swing set.
    \item Estimate the total cost of installing the playground. Then calculate the exact total.
    \item Write equations to represent your calculations for each part.
    \vspace{2cm}
\end{enumerate}
\end{tcolorbox}

\vspace{1em}

% Reflection Box
\begin{tcolorbox}[colframe=black!60, colback=white, 
coltitle=black, colbacktitle=black!15, fonttitle=\bfseries\Large, 
title=Reflection, halign title=center, left=10pt, right=10pt, top=10pt, bottom=100pt]
Reflect on the strategies you used to solve these problems. What patterns did you notice when multiplying or dividing by powers of 10? 
\end{tcolorbox}

\end{document}
